\documentclass[10pt,letterpaper]{article}
\usepackage[margin=0.78in,headheight=17pt]{geometry}
\usepackage[T1]{fontenc}
\usepackage[utf8]{inputenc}
\usepackage{lmodern,microtype,parskip}
\usepackage{amsmath,amssymb}
\usepackage{xcolor,booktabs,longtable,tabularx,array,enumitem}
\usepackage[most]{tcolorbox}
\usepackage{fancyhdr,lastpage,titlesec,hyperref,listings}
\usepackage{tikz}
\usetikzlibrary{arrows.meta,positioning,shapes.geometric}
\definecolor{navy}{HTML}{17324D}
\definecolor{bluegray}{HTML}{43657D}
\definecolor{lightblue}{HTML}{EAF2F8}
\definecolor{linegray}{HTML}{CBD5DF}
\definecolor{softgray}{HTML}{F3F5F7}
\definecolor{accent}{HTML}{8E3B46}
\hypersetup{colorlinks=true,linkcolor=navy,urlcolor=bluegray,pdfborder={0 0 0}}
\urlstyle{same}
\titleformat{\section}{\Large\bfseries\color{navy}}{\thesection}{0.6em}{}
\titleformat{\subsection}{\large\bfseries\color{bluegray}}{\thesubsection}{0.5em}{}
\setlist[itemize]{leftmargin=*,itemsep=3pt,topsep=3pt}
\setlist[enumerate]{leftmargin=*,itemsep=3pt,topsep=3pt}
\setlength{\emergencystretch}{3em}
\renewcommand{\arraystretch}{1.22}
\setlength{\tabcolsep}{4pt}
\pagestyle{fancy}\fancyhf{}
\lhead{Secure Programming \& SAST Engineering}
\rhead{Engineering Documentation Suite}
\cfoot{\thepage\ / \pageref{LastPage}}
\newcolumntype{Y}{>{\raggedright\arraybackslash}X}
\newcommand{\checkbox}{\(\square\)}
\newcommand{\blankline}{\noindent\rule{\linewidth}{0.35pt}\par\vspace{4pt}}
\newcommand{\writearea}[1]{\par\begin{tcolorbox}[colback=white,colframe=linegray,boxrule=0.5pt,height=#1,sharp corners]\mbox{}\end{tcolorbox}\par}
\newtcolorbox{focusbox}[1]{colback=lightblue,colframe=bluegray,title=\textbf{#1},fonttitle=\bfseries,breakable}
\newtcolorbox{warnbox}[1]{colback=softgray,colframe=accent,title=\textbf{#1},fonttitle=\bfseries,breakable}
\lstset{basicstyle=\ttfamily\footnotesize,breaklines=true,frame=single,rulecolor=\color{linegray},backgroundcolor=\color{softgray},showstringspaces=false,columns=fullflexible,keepspaces=true,upquote=true}

\begin{document}
\begin{titlepage}
\vspace*{1.0in}
{\color{bluegray}\Large Chapter study guide}\\[1.0em]
{\color{navy}\Huge\bfseries Chapter 12 Study Guide: Privileged Programs}\\[0.7em]
{\Large Book foundations, guided review, laboratory and modern extension}\\[2.0em]
\rule{\linewidth}{1.0pt}
\vfill
\textbf{Source discipline:} Book-derived foundations are identified separately from current engineering extensions.\\
\textbf{Edition:} 29 October 2026
\end{titlepage}
\tableofcontents
\clearpage
\begin{focusbox}{Book basis and scope}
Chapter 12: Privileged Programs. The 2007 textbook develops the topics listed here; all scenarios and review activities in this guide are independently written teaching exercises, not reproductions of the book.\end{focusbox}
\section{Learning outcomes}
\begin{itemize}
\item Identify privileges held at process and function boundaries.
\item Avoid time-of-check/time-of-use assumptions.
\item Use secure temporary file APIs and argument-vector execution.
\end{itemize}

\section{Concept map: text-derived foundations}
Least privilege, filesystem restrictions, races, temporary files, command injection and inherited file descriptors.
\begin{itemize}
\item More privilege magnifies the consequence of ordinary bugs.
\item File permissions and symlink behavior require careful threat models.
\item External command invocation must not involve attacker-controlled shell text.
\end{itemize}

\section{Reading protocol}
\begin{enumerate}
\item Before reading: define three terms from the chapter title and write one question that a static analyzer might answer.
\item While reading: label each example as threat, faulty assumption, security invariant, corrective pattern, or tool limitation.
\item After reading: draw one data/control-flow diagram and explain what must be established to call a finding confirmed.
\item Record book-specific tooling details separately from generally applicable technical principles.
\end{enumerate}

\section{Application exercise}
\textbf{Scenario.} Compare a temporary-file race design with an atomic secure-create API; document the filesystem threat model.

\textbf{Required artifacts.}\begin{itemize}
\item Scope and trusted versus untrusted inputs; exact code revision and language/runtime.
\item One hypothesis, a rule or manual review path, expected counterexample, and observed result.
\item Documented disposition (true positive, false positive, needs investigation, or risk accepted) and reasoning.
\item A test or bounded argument supporting remediation; explicit limitations and reviewer sign-off.
\end{itemize}

\section{Modern engineering extension (not a claim about the book)}
Add workload identity, container privileges, Linux capabilities and cloud IAM conditions.

\begin{itemize}
\item Map the relevant weakness to an authoritative CWE identifier and cite any applicable current secure coding standard.
\item Treat a current scanner version, command syntax and licensing model as dependencies to verify before implementation.
\item Use immutable commit identifiers, bounded scope and minimal privileges for any repository scan.
\end{itemize}

\section{Knowledge checks}
\begin{enumerate}
\item Which assumption in this chapter can an attacker deliberately violate?
\item Which program representation and analysis sensitivity would improve detection?
\item What realistic false-positive and false-negative mechanisms remain?
\item How would the same security property be verified by a unit test and a manual reviewer?
\item What evidence would you require before waiving the related finding?
\end{enumerate}

\section{Assessment rubric}
\begin{tabularx}{\linewidth}{@{}p{0.22\linewidth}Yp{0.16\linewidth}@{}}\toprule
Criterion & Passing evidence & Points\\\midrule
Conceptual accuracy & Correct threat, failure mechanism and invariant & 0--4\\
Traceability & Links from code location to assertion and finding & 0--4\\
Reproducibility & Tool version, command and sample inputs recorded & 0--4\\
Engineering judgment & Remediation, trade-offs and residual risk explained & 0--4\\
Communication & Concise decision record with no unsupported claims & 0--4\\\bottomrule
\end{tabularx}\par
\textbf{Gate:} 16/20 total and at least 2/4 in every criterion.
\section{Companion suite artifacts}
Use the vulnerability triage, scan baseline and lab evidence workbooks; the program scorecard is for cumulative cohort evaluation. The technical reference defines terms and a standard review lifecycle.
\end{document}
